% 讨论
\section{讨论}
\label{sec:discussion}

\subsection{适用范围与部署模式}

本文方法面向异步带外检测，而非内联阻断。系统在滑动图缓冲区上批量处理审计事件：边级种子识别器先筛出值得调查的事件，链提取和序列模型仅在种子触发时运行。该设计适合持续数小时或数天的 APT 调查场景，牺牲毫秒级阻断能力以换取较高吞吐和可解释的链级输出。

本文评估采用逐日批处理，以便与既有 PIDS 工作公平对比。部署时同一流水线可改为微批流式执行；只要审计事件以时间顺序进入图缓冲区，链提取和 eCDF 校准机制无需改变。

\subsection{局限性}

\textbf{$\branchCoherence$ 不是恶意性分数。}它衡量链是否具有审计支撑的结构一致性，而不是链是否恶意。攻击链和良性链都可能高度一致；恶意性由序列模型的异常分数 $A(\chain)$ 判断。我们通过 \S\ref{sec:metric} 的负控制和敏感性实验验证它确实捕获结构属性，但不把它解释为校准概率。

\textbf{召回受种子识别器约束。}链提取由种子边触发。若攻击步骤完全未被种子识别器浮现，后续链提取无法恢复该步骤。因此端到端召回可分解为种子召回与给定种子覆盖后的链层召回。本文显式报告该分解，以避免把种子识别器能力误归因给链提取器。

\textbf{路径化表示限制。}当前评分对象是以种子为中心的线性化路径，合并后的多个告警链才形成调查图。这适合稀疏、多阶段 APT，但对勒索软件、蠕虫或大规模扫描等分支型攻击不够自然。原生 DAG 评分是直接扩展方向。

\textbf{审计日志覆盖不足的通信通道。}共享内存、部分 \texttt{mmap} IPC 和信号 IPC 在默认审计日志中缺少稳定依赖语义。若攻击主要依赖这些通道，本文方法可能漏检；处理这些通道需要更强的审计 instrumentation 或额外运行时监控。

\textbf{良性训练假设。}本文假设训练数据无攻击污染，这是溯源异常检测中的常见设定。真实部署中可结合鲁棒统计、粗粒度污染过滤或人工确认的良性窗口来放松该假设。

\textbf{跨 engagement 而非完全跨平台。}DARPA E3 与 E5 来自同一计划，虽然场景不同，但审计基础设施相近。因此本文结果只能说明跨 engagement 适配能力；跨 OS、跨审计框架（如 Linux auditd 到 Windows ETW）的泛化仍需进一步验证。

\subsection{崩溃诱导的取证盲区：召回评估的数据上界}
\label{sec:forensic-ceiling}

在核对 CADETS E3 攻击日的原始数据时，我们发现了一个对召回评估有直接影响、且此前工作未曾报告的现象。我们不依赖数据集自带的分片边界，而是对 \texttt{timestampNanos} 做二分查找以精确定位攻击窗口在原始流中的字节偏移。据此我们观测到 04-06 12:08:59--14:01:12 ET 之间存在一段约 1 小时 52 分钟的\emph{物理数据缺口}：该区间内没有任何审计记录。这段缺口与运营事件日志中独立记录的一次故障吻合——``CADETS 目标系统崩溃，原因不明，约 12:20--14:00 ET 停止发布''。

关键在于，该攻击的最后一个活动窗口（约 12:03:50--14:01:32 ET）恰好横跨这段缺口。对照 ground-truth 叙事的末三步——失败的 \texttt{netrecon} 二次探测、\texttt{libdrakon} 载荷投递、以及针对 \texttt{sshd} 的注入尝试——我们在全部约 1330 万条原始记录中做逐字节穷举检索（包括这几步涉及的具体网络地址），命中数为零。也就是说，这三步并非被我们的读取器或种子识别器漏掉，而是\emph{在数据层面根本不存在}。合理的解释是：攻击自身的失稳动作（内存注入尝试）导致主机崩溃，而崩溃连带杀死了审计进程本身，从而摧毁了该攻击后续步骤的溯源证据。无论攻击者是否有意为之，这在客观上构成一种\emph{攻击自毁式的反取证副作用}。

这一现象为召回评估设定了一个与算法无关的数据上界：任何基于溯源图的检测器，无论链提取能力多强，都无法恢复审计日志中不存在的边。因此在该攻击上可达到的召回上界由数据完整性决定（可恢复约 8/11 步），而不仅由种子召回与链层召回决定。据此我们把 \S\ref{sec:evaluation} 的端到端召回分解进一步扩展为
\[
\mathrm{Recall}_{\text{end-to-end}}
= \underbrace{\rho_{\text{data}}}_{\text{数据可恢复性上界}}
\times \mathrm{Recall}_{\text{seed}}
\times \mathrm{Recall}_{\text{chain}\mid\text{seed-covered}},
\]
其中 $\rho_{\text{data}}$ 为落在有效审计覆盖区间内的 ground-truth 步骤占比。我们在报告 CADETS E3 结果时显式标注这三步为``数据层面不可恢复''，仅在可恢复子集上评估链提取质量，以免把数据缺口误算作方法的漏检。

\subsection{未来工作}

未来工作包括四个方向。第一，将当前批处理实现扩展为实时审计日志流处理。第二，设计依赖因果感知的种子识别器，减少对通用 TGN 重构损失的依赖。第三，将 $\branchCoherence$ 从线性路径扩展到原生 DAG 评分，以更好处理分支型攻击。第四，研究跨主机溯源整合；此时依赖边不再完全由单机内核中介，依赖语义也必须相应扩展。
